\section{Introduction}
\label{sec:introduction}

EIP-7702 changes the execution model of an externally owned account (EOA). An authorization tuple names a delegate address; when the authorization is processed, the authority account receives a delegation designator, and calls to that account execute the target's code in the authority's context~\cite{eip7702}. This enables useful account features such as batching, sponsorship, and restricted subkeys without migrating assets to a new address. It also creates a consequential trust decision: code selected by an authorization can act from the signer's established account context. A malicious or inadequately constrained delegate can therefore transfer assets, issue external calls, or exercise privileges that other contracts associate with the EOA~\cite{eip7702,qi2025eip7702phishing,huang2026darkside}.

The defensive decision should consequently precede authorization. Once a user approves hostile code, the harmful behavior may occur in the same interaction or after assets arrive. A newly deployed delegate may have no prior transactions, victims, reputation, or verified source. Transaction-graph and address-reputation systems are valuable after activity accumulates, but those signals may be unavailable to an early signer. Runtime bytecode, by contrast, is a directly inspectable artifact that an integration can obtain at screening time. The question in this work is deliberately narrower than complete wallet protection: can a local model map that bytecode to a useful warning score under a clearly defined label and evaluation boundary?

This setting presents two evaluation hazards. The apparent input may be wrong: EIP-7702 stores the 23-byte designator $\mathtt{0xef0100}\,\|\,\textit{address}$ in the authority account, but the program to inspect is the designated target's runtime bytecode. Treating designators as if they were runtimes lets a bytecode learner consume addresses rather than programs. Exact bytecode also cannot support contradictory class labels without external context. We therefore use an outcome-blind task-alignment policy that removes unresolved designators, retains only safely recovered runtimes, and quarantines complete cross-class exact-bytecode groups without relabeling.

Second, independently sampled rows need not represent independent programs. Contract factories, recompilations, and redeployments produce exact or near-exact relatives. If relatives cross a random train/test boundary, a model can score variants of known code rather than generalize to a new family. We preserve frozen similarity families and outer folds. On this corpus, family-grouped testing controls related-bytecode leakage and provides a more demanding generalization estimate: the same scorer is materially more optimistic under a seeded random diagnostic. Family grouping does not prove attacker independence or eliminate every form of dependence; it makes the measured generalization target explicit.

An attacker can deploy or select a transformed bytecode variant without abandoning the intended attack structure. Metadata, embedded addresses, selectors, and unreachable instructions all influence common bytecode representations. Our benchmark therefore applies structure-preserving transformations under our opcode-skeleton checker and keeps family identity fixed across generated variants. The checker establishes a syntactic invariant, not EVM execution equivalence. We additionally separate a cumulative mutation experiment, a compound flooding stress test, and a stricter augmentation protocol so that results from different attacker conditions are not interpreted as one before/after comparison.

AuthGuard-7702 uses a standard gradient-boosted estimator, and the estimator architecture is not itself a contribution. The AI-tools contribution comes from the implemented bytecode-to-score path together with task-valid input construction, dependence-aware testing, and a reproducible robustness procedure. A conventional model is scientifically useful here because it exposes whether inexpensive learned structure offers signal beyond exact matching and narrow rules, while the protocols reveal when that signal depends on family overlap or fails under transformation. The artifact is an evaluation-grade bytecode-scoring prototype rather than a standalone wallet, authorization parser, RPC service, warning interface, or user study.

Our contributions are:
\begin{itemize}
  \item We implement a bytecode-only, decompiler-free risk-scoring prototype for screening delegate runtime code at the EIP-7702 pre-authorization decision point. Its output is a local score and frozen-threshold warning for an external integration, not a claim of complete wallet protection.
  \item We define an outcome-blind task-alignment and frozen family-grouped evaluation procedure that removes bare designators and conflicting exact-bytecode labels. On the task-aligned dataset, AuthGuard reaches $0.881\pm0.028$ family-grouped AUPRC versus $0.975\pm0.012$ in the seeded random diagnostic, demonstrating optimism from random evaluation on this corpus.
  \item We construct a checker-defined mutation benchmark and a family-disjoint, source-balanced augmentation procedure. For held-out pure-M0 F200, augmentation changes fold-mean recall/AUPRC/FPR from $0.484/0.561/0.217$ to $0.727/0.758/0.174$, with a family-clustered pooled recall difference of $+0.253$ (95\% CI $[0.144,0.379]$), while residual false positives and the compound M3-plus-F200 case remain unresolved.
\end{itemize}
